% ── Sektorové mapy hodinových nákladů práce (PPS/h) ──────────────────────────

Kombinovaný mapový panel (obr.~\ref{fig:sector_wages_map_combined}) porovnává hodinové náklady práce v~paritě kupní síly ve čtyřech klíčových odvětvích \ac{NACE} napříč EU27.
Společná barevná škála umožňuje přímé porovnání absolutní úrovně mezi odvětvími i~zeměmi.
ČR se ve všech čtyřech panelech řadí k~zemím s~nižšími náklady práce, přestože se její produktivita v~řadě odvětví blíží středu EU27.

\begin{description}
  \item[Průmysl (C)]
    ČR navzdory své silné průmyslové základně --- výrobní odvětví tvoří přes \SI{25}{\percent} \ac{HDP}, nejvíce v~\ac{EU} --- dosahuje mzdové úrovně výrazně pod svými produktivitními protihodnotami v~AT nebo DE.
    Mzdová mezera v~průmyslu je emblematickým příkladem fenoménu „low-wage manufacturing periphery": výrobní kapacity přemístěné ze Západu do ČR přinesly zaměstnanost, ale ne přesun mzdové úrovně.
    Tento pattern je udržitelný jen do doby, dokud výrobní náklady v~ČR zůstávají relativně příznivé --- s~rostoucí produktivitou, odlivem pracovní síly a~konkurencí o~pracovníky (srov.~obr.~\ref{fig:emigration_cz_timeline}) se diferenciál snižuje a~zaměstnavatelé budou nuceni ke mzdové konvergenci.
    Proaktivní přístup --- tj.~sektorové \ac{KS} s~pravidelnou indexací průmyslových tarifů --- je z~pohledu pracovníků preferabilní oproti čekání na přirozený tržní tlak.

  \item[Obchod (G)]
    ČR patří mezi státy s~nejnižší mzdovou úrovní v~tomto odvětví --- paradoxně v~situaci, kdy je tento sektor největším zaměstnavatelem v~absolutním počtu osob.
    Maloobchod a~velkoobchod jsou v~ČR charakteristické téměř úplnou absencí \ac{KS} s~plošnou platností: velké maloobchodní řetězce (většinou zahraniční vlastnictví) vyjednávají individuálně nebo mají podnikové \ac{KS} s~minimálními tarify, zatímco tisíce menších zaměstnavatelů nejsou vázány žádnou kolektivní smlouvou.
    Odvětvová \ac{KS} s~erga omnes v~sektoru G by vytvořila jednotný mzdový standard pro všechny maloobchodní zaměstnavatele --- eliminovala by mzdový dumping a~umožnila by lepším zaměstnavatelům soutěžit o~pracovníky na bázi nadstandardu.

  \item[ICT (J)]
    ČR dosahuje relativně lepší pozice než v~ostatních sledovaných odvětvích --- Praha je uznávaným ICT centrem.
    Nicméně i~v~ICT je pozice ČR za AT nebo DE i~po kontrole cenové hladiny: přepočteno na \ac{PPS} je průměrná ICT mzda v~ČR přibližně o~\SIrange{15}{20}{\percent} nižší než v~Německu.
    Paradoxem ICT je, že jde o~sektor s~nejsilnější individuální vyjednávací pozicí --- přesto i~zde absence odvětvové \ac{KS} umožňuje zaměstnavatelům udržet mzdy pod potenciálem, zejména pro méně zkušené pracovníky bez přímé mezinárodní alternativy.
    Odvětvová \ac{KS} v~ICT sektoru by nemusela omezovat tržní prémie pro seniorní odborníky, ale zajistila by minimální tarifní ochranu pro juniorní a~mid-level pracovníky.

  \item[Finance (K)]
    ČR se řadí do středního pole: Pražské finanční centrum vytváří nadprůměrné odměny ve srovnání s~celonárodním průměrem, avšak na mezinárodní úrovni zůstávají pod standardem AT, NL nebo LU.
    Finanční sektor je v~ČR charakteristický oligopolní strukturou se silným podílem zahraničních bank --- zaměstnavatelé v~tomto sektoru mají o~poznání vyšší ziskovost než v~jiných odvětvích, přesto \ac{KS} pokrytí je ve finančním sektoru velmi nízké.
    Mzdová prémie finančního sektoru nad celostátním průměrem je v~ČR nižší než v~AT nebo CH, přestože produktivita na zaměstnance v~bankovnictví je komparabilní nebo vyšší --- přebytek jde primárně do dividendových výnosů zahraničních vlastníků.
    Odvětvová \ac{KS} ve finančním sektoru by mohla zachytit část ziskového přebytku ve prospěch zaměstnanců.
\end{description}

Společným jmenovatelem všech čtyř odvětví je institucionální deficit: absence koordinovaného mzdového vyjednávání umožňuje zaměstnavatelům zachovat mzdovou úroveň pod produktivitním potenciálem ekonomiky.
Zavedení mechanismu erga omnes by tento deficit adresovalo plošně, bez nutnosti sektorově specifické legislativy.

% ── Sektorové mapy hodinových nákladů práce (PPS/h) ──────────────────────────

Kombinovaný mapový panel (obr.~\ref{fig:sector_wages_map_combined}) porovnává hodinové náklady práce v~paritě kupní síly ve čtyřech klíčových odvětvích \ac{NACE} napříč EU27.
Společná barevná škála umožňuje přímé porovnání absolutní úrovně mezi odvětvími i~zeměmi.
ČR se ve všech čtyřech panelech řadí k~zemím s~nižšími náklady práce, přestože se její produktivita v~řadě odvětví blíží středu EU27.

\begin{description}
  \item[Průmysl (C)]
    ČR navzdory své silné průmyslové základně --- výrobní odvětví tvoří přes \SI{25}{\percent} \ac{HDP}, nejvíce v~\ac{EU} --- dosahuje mzdové úrovně výrazně pod svými produktivitními protihodnotami v~AT nebo DE.
    Mzdová mezera v~průmyslu je emblematickým příkladem fenoménu „low-wage manufacturing periphery": výrobní kapacity přemístěné ze Západu do ČR přinesly zaměstnanost, ale ne přesun mzdové úrovně.
    Tento pattern je udržitelný jen do doby, dokud výrobní náklady v~ČR zůstávají relativně příznivé --- s~rostoucí produktivitou, odlivem pracovní síly a~konkurencí o~pracovníky (srov.~obr.~\ref{fig:emigration_cz_timeline}) se diferenciál snižuje a~zaměstnavatelé budou nuceni ke mzdové konvergenci.
    Proaktivní přístup --- tj.~sektorové \ac{KS} s~pravidelnou indexací průmyslových tarifů --- je z~pohledu pracovníků preferabilní oproti čekání na přirozený tržní tlak.

  \item[Obchod (G)]
    ČR patří mezi státy s~nejnižší mzdovou úrovní v~tomto odvětví --- paradoxně v~situaci, kdy je tento sektor největším zaměstnavatelem v~absolutním počtu osob.
    Maloobchod a~velkoobchod jsou v~ČR charakteristické téměř úplnou absencí \ac{KS} s~plošnou platností: velké maloobchodní řetězce (většinou zahraniční vlastnictví) vyjednávají individuálně nebo mají podnikové \ac{KS} s~minimálními tarify, zatímco tisíce menších zaměstnavatelů nejsou vázány žádnou kolektivní smlouvou.
    Odvětvová \ac{KS} s~erga omnes v~sektoru G by vytvořila jednotný mzdový standard pro všechny maloobchodní zaměstnavatele --- eliminovala by mzdový dumping a~umožnila by lepším zaměstnavatelům soutěžit o~pracovníky na bázi nadstandardu.

  \item[ICT (J)]
    ČR dosahuje relativně lepší pozice než v~ostatních sledovaných odvětvích --- Praha je uznávaným ICT centrem.
    Nicméně i~v~ICT je pozice ČR za AT nebo DE i~po kontrole cenové hladiny: přepočteno na \ac{PPS} je průměrná ICT mzda v~ČR přibližně o~\SIrange{15}{20}{\percent} nižší než v~Německu.
    Paradoxem ICT je, že jde o~sektor s~nejsilnější individuální vyjednávací pozicí --- přesto i~zde absence odvětvové \ac{KS} umožňuje zaměstnavatelům udržet mzdy pod potenciálem, zejména pro méně zkušené pracovníky bez přímé mezinárodní alternativy.
    Odvětvová \ac{KS} v~ICT sektoru by nemusela omezovat tržní prémie pro seniorní odborníky, ale zajistila by minimální tarifní ochranu pro juniorní a~mid-level pracovníky.

  \item[Finance (K)]
    ČR se řadí do středního pole: Pražské finanční centrum vytváří nadprůměrné odměny ve srovnání s~celonárodním průměrem, avšak na mezinárodní úrovni zůstávají pod standardem AT, NL nebo LU.
    Finanční sektor je v~ČR charakteristický oligopolní strukturou se silným podílem zahraničních bank --- zaměstnavatelé v~tomto sektoru mají o~poznání vyšší ziskovost než v~jiných odvětvích, přesto \ac{KS} pokrytí je ve finančním sektoru velmi nízké.
    Mzdová prémie finančního sektoru nad celostátním průměrem je v~ČR nižší než v~AT nebo CH, přestože produktivita na zaměstnance v~bankovnictví je komparabilní nebo vyšší --- přebytek jde primárně do dividendových výnosů zahraničních vlastníků.
    Odvětvová \ac{KS} ve finančním sektoru by mohla zachytit část ziskového přebytku ve prospěch zaměstnanců.
\end{description}

Společným jmenovatelem všech čtyř odvětví je institucionální deficit: absence koordinovaného mzdového vyjednávání umožňuje zaměstnavatelům zachovat mzdovou úroveň pod produktivitním potenciálem ekonomiky.
Zavedení mechanismu erga omnes by tento deficit adresovalo plošně, bez nutnosti sektorově specifické legislativy.

% ── Sektorové mapy hodinových nákladů práce (PPS/h) ──────────────────────────

Kombinovaný mapový panel (obr.~\ref{fig:sector_wages_map_combined}) porovnává hodinové náklady práce v~paritě kupní síly ve čtyřech klíčových odvětvích \ac{NACE} napříč EU27.
Společná barevná škála umožňuje přímé porovnání absolutní úrovně mezi odvětvími i~zeměmi.
ČR se ve všech čtyřech panelech řadí k~zemím s~nižšími náklady práce, přestože se její produktivita v~řadě odvětví blíží středu EU27.

\begin{description}
  \item[Průmysl (C)]
    ČR navzdory své silné průmyslové základně --- výrobní odvětví tvoří přes \SI{25}{\percent} \ac{HDP}, nejvíce v~\ac{EU} --- dosahuje mzdové úrovně výrazně pod svými produktivitními protihodnotami v~AT nebo DE.
    Mzdová mezera v~průmyslu je emblematickým příkladem fenoménu „low-wage manufacturing periphery": výrobní kapacity přemístěné ze Západu do ČR přinesly zaměstnanost, ale ne přesun mzdové úrovně.
    Tento pattern je udržitelný jen do doby, dokud výrobní náklady v~ČR zůstávají relativně příznivé --- s~rostoucí produktivitou, odlivem pracovní síly a~konkurencí o~pracovníky (srov.~obr.~\ref{fig:emigration_cz_timeline}) se diferenciál snižuje a~zaměstnavatelé budou nuceni ke mzdové konvergenci.
    Proaktivní přístup --- tj.~sektorové \ac{KS} s~pravidelnou indexací průmyslových tarifů --- je z~pohledu pracovníků preferabilní oproti čekání na přirozený tržní tlak.

  \item[Obchod (G)]
    ČR patří mezi státy s~nejnižší mzdovou úrovní v~tomto odvětví --- paradoxně v~situaci, kdy je tento sektor největším zaměstnavatelem v~absolutním počtu osob.
    Maloobchod a~velkoobchod jsou v~ČR charakteristické téměř úplnou absencí \ac{KS} s~plošnou platností: velké maloobchodní řetězce (většinou zahraniční vlastnictví) vyjednávají individuálně nebo mají podnikové \ac{KS} s~minimálními tarify, zatímco tisíce menších zaměstnavatelů nejsou vázány žádnou kolektivní smlouvou.
    Odvětvová \ac{KS} s~erga omnes v~sektoru G by vytvořila jednotný mzdový standard pro všechny maloobchodní zaměstnavatele --- eliminovala by mzdový dumping a~umožnila by lepším zaměstnavatelům soutěžit o~pracovníky na bázi nadstandardu.

  \item[ICT (J)]
    ČR dosahuje relativně lepší pozice než v~ostatních sledovaných odvětvích --- Praha je uznávaným ICT centrem.
    Nicméně i~v~ICT je pozice ČR za AT nebo DE i~po kontrole cenové hladiny: přepočteno na \ac{PPS} je průměrná ICT mzda v~ČR přibližně o~\SIrange{15}{20}{\percent} nižší než v~Německu.
    Paradoxem ICT je, že jde o~sektor s~nejsilnější individuální vyjednávací pozicí --- přesto i~zde absence odvětvové \ac{KS} umožňuje zaměstnavatelům udržet mzdy pod potenciálem, zejména pro méně zkušené pracovníky bez přímé mezinárodní alternativy.
    Odvětvová \ac{KS} v~ICT sektoru by nemusela omezovat tržní prémie pro seniorní odborníky, ale zajistila by minimální tarifní ochranu pro juniorní a~mid-level pracovníky.

  \item[Finance (K)]
    ČR se řadí do středního pole: Pražské finanční centrum vytváří nadprůměrné odměny ve srovnání s~celonárodním průměrem, avšak na mezinárodní úrovni zůstávají pod standardem AT, NL nebo LU.
    Finanční sektor je v~ČR charakteristický oligopolní strukturou se silným podílem zahraničních bank --- zaměstnavatelé v~tomto sektoru mají o~poznání vyšší ziskovost než v~jiných odvětvích, přesto \ac{KS} pokrytí je ve finančním sektoru velmi nízké.
    Mzdová prémie finančního sektoru nad celostátním průměrem je v~ČR nižší než v~AT nebo CH, přestože produktivita na zaměstnance v~bankovnictví je komparabilní nebo vyšší --- přebytek jde primárně do dividendových výnosů zahraničních vlastníků.
    Odvětvová \ac{KS} ve finančním sektoru by mohla zachytit část ziskového přebytku ve prospěch zaměstnanců.
\end{description}

Společným jmenovatelem všech čtyř odvětví je institucionální deficit: absence koordinovaného mzdového vyjednávání umožňuje zaměstnavatelům zachovat mzdovou úroveň pod produktivitním potenciálem ekonomiky.
Zavedení mechanismu erga omnes by tento deficit adresovalo plošně, bez nutnosti sektorově specifické legislativy.

% ── Sektorové mapy hodinových nákladů práce (PPS/h) ──────────────────────────

Kombinovaný mapový panel (obr.~\ref{fig:sector_wages_map_combined}) porovnává hodinové náklady práce v~paritě kupní síly ve čtyřech klíčových odvětvích \ac{NACE} napříč EU27.
Společná barevná škála umožňuje přímé porovnání absolutní úrovně mezi odvětvími i~zeměmi.
ČR se ve všech čtyřech panelech řadí k~zemím s~nižšími náklady práce, přestože se její produktivita v~řadě odvětví blíží středu EU27.

\begin{description}
  \item[Průmysl (C)]
    ČR navzdory své silné průmyslové základně --- výrobní odvětví tvoří přes \SI{25}{\percent} \ac{HDP}, nejvíce v~\ac{EU} --- dosahuje mzdové úrovně výrazně pod svými produktivitními protihodnotami v~AT nebo DE.
    Mzdová mezera v~průmyslu je emblematickým příkladem fenoménu „low-wage manufacturing periphery": výrobní kapacity přemístěné ze Západu do ČR přinesly zaměstnanost, ale ne přesun mzdové úrovně.
    Tento pattern je udržitelný jen do doby, dokud výrobní náklady v~ČR zůstávají relativně příznivé --- s~rostoucí produktivitou, odlivem pracovní síly a~konkurencí o~pracovníky (srov.~obr.~\ref{fig:emigration_cz_timeline}) se diferenciál snižuje a~zaměstnavatelé budou nuceni ke mzdové konvergenci.
    Proaktivní přístup --- tj.~sektorové \ac{KS} s~pravidelnou indexací průmyslových tarifů --- je z~pohledu pracovníků preferabilní oproti čekání na přirozený tržní tlak.

  \item[Obchod (G)]
    ČR patří mezi státy s~nejnižší mzdovou úrovní v~tomto odvětví --- paradoxně v~situaci, kdy je tento sektor největším zaměstnavatelem v~absolutním počtu osob.
    Maloobchod a~velkoobchod jsou v~ČR charakteristické téměř úplnou absencí \ac{KS} s~plošnou platností: velké maloobchodní řetězce (většinou zahraniční vlastnictví) vyjednávají individuálně nebo mají podnikové \ac{KS} s~minimálními tarify, zatímco tisíce menších zaměstnavatelů nejsou vázány žádnou kolektivní smlouvou.
    Odvětvová \ac{KS} s~erga omnes v~sektoru G by vytvořila jednotný mzdový standard pro všechny maloobchodní zaměstnavatele --- eliminovala by mzdový dumping a~umožnila by lepším zaměstnavatelům soutěžit o~pracovníky na bázi nadstandardu.

  \item[ICT (J)]
    ČR dosahuje relativně lepší pozice než v~ostatních sledovaných odvětvích --- Praha je uznávaným ICT centrem.
    Nicméně i~v~ICT je pozice ČR za AT nebo DE i~po kontrole cenové hladiny: přepočteno na \ac{PPS} je průměrná ICT mzda v~ČR přibližně o~\SIrange{15}{20}{\percent} nižší než v~Německu.
    Paradoxem ICT je, že jde o~sektor s~nejsilnější individuální vyjednávací pozicí --- přesto i~zde absence odvětvové \ac{KS} umožňuje zaměstnavatelům udržet mzdy pod potenciálem, zejména pro méně zkušené pracovníky bez přímé mezinárodní alternativy.
    Odvětvová \ac{KS} v~ICT sektoru by nemusela omezovat tržní prémie pro seniorní odborníky, ale zajistila by minimální tarifní ochranu pro juniorní a~mid-level pracovníky.

  \item[Finance (K)]
    ČR se řadí do středního pole: Pražské finanční centrum vytváří nadprůměrné odměny ve srovnání s~celonárodním průměrem, avšak na mezinárodní úrovni zůstávají pod standardem AT, NL nebo LU.
    Finanční sektor je v~ČR charakteristický oligopolní strukturou se silným podílem zahraničních bank --- zaměstnavatelé v~tomto sektoru mají o~poznání vyšší ziskovost než v~jiných odvětvích, přesto \ac{KS} pokrytí je ve finančním sektoru velmi nízké.
    Mzdová prémie finančního sektoru nad celostátním průměrem je v~ČR nižší než v~AT nebo CH, přestože produktivita na zaměstnance v~bankovnictví je komparabilní nebo vyšší --- přebytek jde primárně do dividendových výnosů zahraničních vlastníků.
    Odvětvová \ac{KS} ve finančním sektoru by mohla zachytit část ziskového přebytku ve prospěch zaměstnanců.
\end{description}

Společným jmenovatelem všech čtyř odvětví je institucionální deficit: absence koordinovaného mzdového vyjednávání umožňuje zaměstnavatelům zachovat mzdovou úroveň pod produktivitním potenciálem ekonomiky.
Zavedení mechanismu erga omnes by tento deficit adresovalo plošně, bez nutnosti sektorově specifické legislativy.

\input{texparts/python/sector_wages_map_combined}



